% Generated draft: requires mathematical and source review.
\documentclass[11pt]{article}
\usepackage[margin=25mm]{geometry}
\usepackage{fontspec}
\setmainfont{Latin Modern Roman}
\usepackage{amsmath,amssymb,mathtools,unicode-math}
\setmathfont{Latin Modern Math}
\usepackage{enumitem,needspace,xurl,hyperref}
\setlength{\parindent}{0pt}
\setlength{\parskip}{5pt}
\setlength{\emergencystretch}{3em}
\newcommand{\N}{\mathbb{N}}
\newcommand{\Z}{\mathbb{Z}}
\newcommand{\Q}{\mathbb{Q}}
\newcommand{\R}{\mathbb{R}}
\newcommand{\C}{\mathbb{C}}
\newcommand{\F}{\mathbb{F}}
\newcommand{\A}{\mathbb{A}}
\newcommand{\PP}{\mathbb{P}}
\newcommand{\E}{\mathbb{E}}
\newcommand{\eps}{\varepsilon}
\newcommand{\dd}{\,\mathrm{d}}
\begin{document}
\section*{UnsolvedMath ID: 800752}
\textbf{Source problem code:} catalogue entry
\subsection{Definitions and mathematical statement}
Write \(\N=\{0,1,2,\ldots\}\). For \(f,g\in\N^{\N}\), put \(f\leq^{*}g\) if \(f(n)\leq g(n)\) for all but finitely many \(n\). The bounding number is
\[
\mathfrak b=\min\bigl\{|B|:B\subseteq\N^{\N},\ (\forall g\in\N^{\N})(\exists f\in B)\ f\not\leq^{*}g\bigr\}.
\]
Here \(\aleph_2\) denotes the second uncountable cardinal.

For a set \(S\), let \([S]^{\leq\aleph_0}\) denote its finite and countably infinite subsets. An ideal \(\mathcal I\subseteq[S]^{\leq\aleph_0}\) contains every finite subset of \(S\) and is closed under taking subsets and finite unions. It is a \(P\)-ideal if every sequence \((A_n)_{n\in\N}\) in \(\mathcal I\) has an \(A\in\mathcal I\) such that \(A_n\setminus A\) is finite for every \(n\). The \emph{\(P\)-ideal dichotomy} (PID) asserts that, for every set \(S\) and every such \(P\)-ideal \(\mathcal I\), at least one of the following holds:
\begin{enumerate}
\item There is an uncountable \(T\subseteq S\) such that \([T]^{\leq\aleph_0}\subseteq\mathcal I\).
\item There are sets \(S_n\subseteq S\), for \(n\in\N\), with \(S=\bigcup_{n\in\N}S_n\), such that \(S_n\cap A\) is finite for every \(n\in\N\) and every \(A\in\mathcal I\).
\end{enumerate}

Let \(\mathbb K\in\{\R,\C\}\) be a scalar field. A Banach space over \(\mathbb K\) is a complete normed vector space over \(\mathbb K\); it is nonseparable if it has no countable norm-dense subset. Write \(X^*\) for its continuous \(\mathbb K\)-linear dual. An uncountable biorthogonal system in \(X\) is a family \((x_\alpha,x_\alpha^*)_{\alpha\in A}\), where \(A\) is uncountable, \(x_\alpha\in X\), \(x_\alpha^*\in X^*\), and
\[
x_\alpha^*(x_\beta)=\begin{cases}1,&\alpha=\beta,\\0,&\alpha\ne\beta\end{cases}
\qquad(\alpha,\beta\in A).
\]
Let \(\mathcal B_{\mathbb K}\) be the assertion that every nonseparable Banach space over \(\mathbb K\) admits such a system.

The supplied Problem 16 asks: assuming PID, are the following assertions equivalent?
\[
\mathcal B_{\mathbb K}\qquad\Longleftrightarrow\qquad\mathfrak b=\aleph_2.
\]
The supplied source record does not specify whether its Banach spaces are real or complex. The parameter \(\mathbb K\) makes this unresolved convention explicit; it does not assert that the source separately asks the question for both fields.
\subsection{Short English statement}
Assuming the \(P\)-ideal dichotomy, is the assertion that every nonseparable Banach space has an uncountable biorthogonal system equivalent to \(\mathfrak b=\aleph_2\)? The supplied record leaves the scalar field unspecified and labels the problem open; that status has not been independently verified.
\subsection{Sources}
\begin{itemize}
\item Supplied local UnsolvedMath registry record: permanent ID \(800752\), source code \texttt{catalogue entry}, title \emph{Banach Spaces and Ramsey Theory --- Problem 16}. Attribution in the record: Dodos and Lopez-Abad, proposed year 2010. The registry supplies the status \texttt{open}.
\item Dodos and Lopez-Abad, \emph{Banach Spaces and Ramsey Theory: Some Open Problems} (2010), Problem 16. Supplied source URL: \url{https://arxiv.org/abs/1006.2668}. The background records access on 2026-10-16 and extraction by \texttt{arxiv\_latex}; these are supplied provenance statements, not independently verified access or source checks.
\end{itemize}
\end{document}
